\section{유한성 조건}
\label{section:I.6-ko}

\subsection{뇌터 준스킴과 국소 뇌터 준스킴}
\label{subsection:I.6.1-ko}

\begin{definition}[6.1.1]
\label{I.6.1.1-ko}
준스킴 $X$가 환이 뇌터인 아핀 열린집합 $V_\alpha$들의 유한 합집합
(각각 합집합)이면 $X$를 \emph{뇌터}(각각 \emph{국소 뇌터})라 한다.
여기서 각 환은 $V_\alpha$ 위에 유도된 스킴의 환을 뜻한다.
\end{definition}

$X$가 국소 뇌터이면 구조층 $\mathscr{O}_X$가
\emph{연접인 환의 층}이라는 것은 문제가 국소적임과
(\hyperref[I.1.5.2-ko]{1.5.2})에서 곧바로 따른다.
$\mathscr{O}_X$-가군 $\mathscr{F}$가 \emph{연접}이면,
$\mathscr{F}$의 \emph{모든 준연접 $\mathscr{O}_X$-부분가군}
\oldpage[I]{141}
\emph{(각각 모든 준연접 $\mathscr{O}_X$-몫가군)}은
\emph{연접이다}.
실제로 이 문제도 국소적이며, 뇌터 환 위의 유한 생성 가군의 부분가군
(각각 몫가군)이 유한 생성이라는 사실과
(\hyperref[I.1.5.1-ko]{1.5.1}),
(\hyperref[I.1.4.1-ko]{1.4.1}),
(\hyperref[I.1.3.10-ko]{1.3.10})을 적용하면 충분하다. 특히
$\mathscr{O}_X$의 \emph{모든 준연접 아이디얼층은} \emph{연접이다}.

준스킴 $X$가 열린집합 $W_\lambda$들의 유한 합집합(각각 합집합)이고,
각 $W_\lambda$ 위에 유도된 준스킴이 뇌터(각각 국소 뇌터)이면 $X$가
뇌터(각각 국소 뇌터)임은 자명하다.

\begin{proposition}[6.1.2]
\label{I.6.1.2-ko}
준스킴 $X$가 뇌터이기 위한 필요충분조건은 $X$가 국소 뇌터이고 그
바탕공간이 준콤팩트인 것이다. 이때 $X$의 바탕공간은 뇌터이다.
\end{proposition}

\begin{proof}
첫째 주장은 정의와 (\hyperref[I.1.1.10-ko]{1.1.10 (ii)})에서 곧바로
따른다. 둘째 주장은 (\hyperref[I.1.1.6-ko]{1.1.6})과 뇌터
부분공간들의 유한 합집합인 모든 공간이 뇌터라는 사실
(\hyperref[0.2.2.3-ko]{0, 2.2.3})에서 따른다.
\end{proof}

\begin{proposition}[6.1.3]
\label{I.6.1.3-ko}
$X$를 환 $A$의 아핀 스킴이라 하자. 다음 조건들은 동치이다.
a) $X$는 뇌터이다. b) $X$는 국소 뇌터이다. c) $A$는 뇌터이다.
\end{proposition}

\begin{proof}
a)와 b)의 동치는 (\hyperref[I.6.1.2-ko]{6.1.2})와 모든 아핀 스킴의
바탕공간이 준콤팩트라는 사실
(\hyperref[I.1.1.10-ko]{1.1.10})에서 따른다. 또한 c)가 a)를
함의함은 자명하다. a)가 c)를 함의함을 보이자. 환이 뇌터인 아핀
열린집합 $V_i$들의 유한 덮개를 $X$에서 잡고, $V_i$ 위에 유도된
준스킴의 환을 $A_i$라 하자. $A$의 아이디얼들의 증가열
$(\mathfrak{a}_n)$을 잡으면 이에 표준적으로 전단사 대응하는
(\hyperref[I.1.3.7-ko]{1.3.7})
$\widetilde{A}\equiv\mathscr{O}_X$의 아이디얼층들의 증가열
$(\widetilde{\mathfrak{a}}_n)$이 존재한다. $(\mathfrak{a}_n)$이
어느 단계 이후 일정함을 보이려면 $(\widetilde{\mathfrak{a}}_n)$이
어느 단계 이후 일정함을
보이면 충분하다. 그런데 $\widetilde{\mathfrak{a}}_n|V_i$는 표준
포함 사상 $V_i\to X$에 의한 $\widetilde{\mathfrak{a}}_n$의
역상이므로 (\hyperref[0.5.1.4-ko]{0, 5.1.4})
$\mathscr{O}_X|V_i$의 준연접 아이디얼층이다. 따라서
$\widetilde{\mathfrak{a}}_n|V_i=\widetilde{\mathfrak{a}}_{ni}$이고,
여기서 $\mathfrak{a}_{ni}$는 $A_i$의 아이디얼이다
(\hyperref[I.1.3.7-ko]{1.3.7}). $A_i$가 뇌터이므로 모든 $i$에
대하여 $(\mathfrak{a}_{ni})$는 어느 단계 이후 일정하며, 이로부터
명제가 따른다.
\end{proof}

앞의 논증은 \emph{$X$가 뇌터 준스킴이면 $\mathscr{O}_X$의 연접
아이디얼층들의 모든 증가열은 어느 단계 이후 일정하다}라는 것도
증명한다.

\begin{proposition}[6.1.4]
\label{I.6.1.4-ko}
뇌터(각각 국소 뇌터) 준스킴의 모든 부분준스킴은 뇌터(각각 국소
뇌터)이다.
\end{proposition}

\begin{proof}
뇌터 준스킴 $X$에 대해서만 증명하면 충분하다. 또한 정의
(\hyperref[I.6.1.1-ko]{6.1.1})에 따라 $X$가 아핀 스킴인 경우로
곧바로 환원된다. $X$의 모든 부분준스킴은 어떤 열린집합 위에 유도된
준스킴의 닫힌 부분준스킴이므로
(\hyperref[I.4.1.3-ko]{4.1.3}), $Y$가 $X$의 닫힌 부분준스킴이거나
열린집합 위에 유도된 부분준스킴인 경우만 생각하면 된다. $Y$가
닫힌 경우는 자명하다. 실제로 $A$가 $X$의 환이면 $Y$는 $A$의 어떤
아이디얼 $\mathfrak{J}$에 대하여 환 $A/\mathfrak{J}$의 아핀
스킴이다(\hyperref[I.4.2.3-ko]{4.2.3}). $A$가 뇌터이므로
(\hyperref[I.6.1.3-ko]{6.1.3}) $A/\mathfrak{J}$도 뇌터이다.

이제 $Y$가 $X$의 열린집합이라고 하자. 바탕공간 $Y$는 뇌터이므로
(\hyperref[I.6.1.2-ko]{6.1.2}) 준콤팩트이고, 따라서
$f_i\in A$인 열린집합 $D(f_i)$들의 유한 합집합이다.
\oldpage[I]{142}
그러므로 $f\in A$에 대하여 $Y=D(f)$인 경우만 증명하면 된다. 이때
$Y$는 환이 $A_f$와 동형인 아핀 스킴이다
(\hyperref[I.1.3.6-ko]{1.3.6}). $A$가 뇌터이므로
(\hyperref[I.6.1.3-ko]{6.1.3}) $A_f$도 뇌터이다.
\end{proof}

\begin{env}[6.1.5]
\label{I.6.1.5-ko}
두 뇌터 $S$-준스킴의 \emph{곱}은 반드시 뇌터인 것은 아니다. 이는
두 뇌터 대수의 텐서곱이 반드시 뇌터 환인 것은 아니므로, 두
준스킴이 아핀인 경우에도 마찬가지이다
(\agref{I.6.3.8-ko}{6.3.8} 참조).
\end{env}

\begin{proposition}[6.1.6]
\label{I.6.1.6-ko}
$X$가 뇌터 준스킴이면 $\mathscr{O}_X$의 멱영근기
$\mathscr{N}_X$는 멱영이다.
\end{proposition}

실제로 $X$를 유한개의 아핀 열린집합 $U_i$로 덮을 수 있고,
$(\mathscr{N}_X|U_i)^{n_i}=0$인 정수 $n_i$가 존재함을 보이면
충분하다. $n$을 $n_i$들의 최댓값이라 하면
$\mathscr{N}_X^n=0$이다. 따라서 $A$가 뇌터 환이고
$X=\operatorname{Spec}(A)$가 아핀인 경우로 환원된다. 이때
(\hyperref[I.5.1.1-ko]{5.1.1})과
(\hyperref[I.1.3.13-ko]{1.3.13})에 따라 $A$의 멱영근기가 멱영이라는
사실(\cite{I-11}, p.~127, cor.~4)을 확인하면 충분하다.

\begin{corollary}[6.1.7]
\label{I.6.1.7-ko}
$X$를 뇌터 준스킴이라 하자. $X$가 아핀 스킴이기 위한
필요충분조건은 $X_{\mathrm{red}}$가 아핀 스킴인 것이다.
\end{corollary}

이는 (\hyperref[I.6.1.6-ko]{6.1.6})과
(\hyperref[I.5.1.10-ko]{5.1.10})에서 따른다.

\begin{lemma}[6.1.8]
\label{I.6.1.8-ko}
$X$를 위상공간, $x$를 $X$의 한 점, $U$를 기약 성분이 유한개뿐인
$x$의 열린 근방이라 하자. 그러면 $x$의 근방 $V$가 존재하여,
$V$에 포함되는 $x$의 모든 열린 근방이 연결이다.
\end{lemma}

$x$를 포함하지 않는 $U$의 기약 성분들을 $U_i$($1\leq i\leq m$)라
하자. 그 합집합의 $U$에서의 여집합은 $U$ 안에서, 따라서 $X$ 안에서
$x$의 열린 근방 $V$이다. 또한 이는 $U$와, $x$를 포함하지 않는
$X$의 기약 성분들의 합집합을 $X$에서 뺀 여집합의 교집합이다
(\hyperref[0.2.1.6-ko]{0, 2.1.6}). 이제 $V$에 포함되는 $x$의 열린
근방 $W$를 잡자. $W$의 기약 성분들은 $W$와 만나는 $U$의 기약
성분들과 $W$의 교집합들이므로
(\hyperref[0.2.1.6-ko]{0, 2.1.6}), 모두 $x$를 포함한다. 이 성분들은
연결이므로 $W$도 연결이다.

\begin{corollary}[6.1.9]
\label{I.6.1.9-ko}
국소 뇌터 위상공간은 국소 연결이다. 따라서 특히 그 연결 성분들은
열린집합이다.
\end{corollary}

\begin{proposition}[6.1.10]
\label{I.6.1.10-ko}
$X$를 국소 뇌터 위상공간이라 하자. 다음 조건들은 동치이다.
\begin{enumerate}
  \item[a)] $X$의 기약 성분들은 열린집합이다.
  \item[b)] $X$의 기약 성분들은 연결 성분들과 일치한다.
  \item[c)] $X$의 연결 성분들은 기약이다.
  \item[d)] $X$의 서로 다른 두 기약 성분은 만나지 않는다.
\end{enumerate}

또한 $X$가 준스킴이면 이 조건들은 다음 조건과도 동치이다.
\begin{enumerate}
  \item[e)] 모든 $x\in X$에 대하여
    $\operatorname{Spec}(\mathscr{O}_x)$는 기약이다. 다시 말해
    $\mathscr{O}_x$의 멱영근기는 소 아이디얼이다.
\end{enumerate}
\end{proposition}

기약공간은 연결이므로 a)는 b)를 함의한다. 실제로 a)에 따라 $X$의
기약 성분들은 동시에 열린집합이고 닫힌집합이다. b)가 c)를 함의함은
자명하다. 역으로 닫힌집합 $F$가
\oldpage[I]{143}
$X$의 연결 성분 $C$를 포함하면서 $C$와 다르면 $F$는 기약일 수
없다. 실제로 이 집합은 연결이 아니므로
$F$ 안에서 동시에 열리고 닫힌 서로소인 두 공집합 아닌 집합의
합집합이며, 이 두 집합은 $X$에서 닫혀 있다. 따라서 c)는 b)를
함의한다. 이로부터 서로 다른 두 연결 성분은 만나지 않으므로 c)가
d)를 함의함이 곧바로 따른다.

여기까지는 $X$가 국소 뇌터라는 사실을 사용하지 않았다. 이제 이
가정을 하고 d)가 a)를 함의함을 보이자. (\hyperref[0.2.1.6-ko]{0,
2.1.6})에 따라 $X$가 뇌터인 경우로 환원할 수 있고, 이때 $X$의 기약
성분은 유한개뿐이다. 이 성분들이 닫혀 있고 서로소이므로 열린집합이다.

끝으로 d)와 e)의 동치는 준스킴 $X$의 바탕공간이 국소 뇌터라는
가정 없이도 성립한다. (\hyperref[0.2.1.6-ko]{0, 2.1.6})에 따라
$X=\operatorname{Spec}(A)$가 아핀인 경우로 환원할 수 있다. $x$가
$X$의 기약 성분 하나에만 포함된다는 것은
$\mathfrak{j}_x$가 $A$의 극소 소 아이디얼 하나만을 포함한다는 뜻이고
(\hyperref[I.1.1.14-ko]{1.1.14}), 이는
$\mathfrak{j}_x\mathscr{O}_x$가 $\mathscr{O}_x$의 극소 소 아이디얼
하나만을 포함한다는 것과 동치이다. 이로부터 결론이 따른다.

\begin{corollary}[6.1.11]
\label{I.6.1.11-ko}
$X$를 국소 뇌터 공간이라 하자. $X$가 기약이기 위한 필요충분조건은
$X$가 연결이고 공집합이 아니며, $X$의 서로 다른 두 기약 성분이
만나지 않는 것이다. $X$가 준스킴이면 마지막 조건은 모든 $x\in X$에
대하여 $\operatorname{Spec}(\mathscr{O}_x)$가 기약이라는 것과
동치이다.
\end{corollary}

마지막 주장은 (\hyperref[I.6.1.10-ko]{6.1.10})에서 보였다. 따라서
첫째 주장의 조건들이 충분함만 증명하면 된다. 그런데
(\hyperref[I.6.1.10-ko]{6.1.10})에 따라 이 조건들은 $X$의 기약
성분들이 연결 성분들과 일치함을 함의하고, $X$가 연결이고 공집합이
아니므로 $X$는 기약이다.

\begin{corollary}[6.1.12]
\label{I.6.1.12-ko}
$X$를 국소 뇌터 준스킴이라 하자. $X$가 정역 준스킴이기 위한
필요충분조건은 $X$가 공집합이 아니고 연결이며 모든 $x\in X$에 대하여
$\mathscr{O}_x$가 정역인 것이다.
\end{corollary}

\begin{proposition}[6.1.13]
\label{I.6.1.13-ko}
$X$를 국소 뇌터 준스킴, $x\in X$를 한 점이라 하자. 국소환
$\mathscr{O}_x$의 멱영근기 $\mathscr{N}_x$가 소 아이디얼이면(각각
$\mathscr{O}_x$가 축소환이면, 정역이면), $x$의 기약인(각각 축소인,
정역인) 열린 근방 $U$가 존재한다.
\end{proposition}

$\mathscr{N}_x$가 소 아이디얼인 경우와 $\mathscr{N}_x=0$인 경우만
생각하면 충분하다. 셋째 가정은 이 두 가정을 함께 말한 것이기
때문이다. $\mathscr{N}_x$가 소 아이디얼이면 $x$는 $X$의 기약 성분
$Y$ 하나에만 속한다(\hyperref[I.6.1.10-ko]{6.1.10}). $x$를 포함하지
않는 $X$의 기약 성분들의 집합은 국소 유한이므로 그 합집합은 닫혀
있다. 따라서 이 합집합의 여집합 $U$는 열린집합이고 $Y$에 포함되므로
기약이다(\hyperref[0.2.1.6-ko]{0, 2.1.6}). $\mathscr{N}_x=0$이면
$x$의 어떤 근방에 속하는 모든 $y$에 대하여 $\mathscr{N}_y=0$이다.
실제로 $\mathscr{N}$은 준연접이고
(\hyperref[I.5.1.1-ko]{5.1.1}), $X$가 국소 뇌터이므로 연접이다.
따라서 결론은 (\hyperref[0.5.2.2-ko]{0, 5.2.2})에서 따른다.

\subsection{아르틴 준스킴}
\label{subsection:I.6.2-ko}

\begin{definition}[6.2.1]
\label{I.6.2.1-ko}
준스킴이 아핀이고 그 환이 아르틴 환이면 그 준스킴을
\emph{아르틴}이라고 한다.
\end{definition}

\oldpage[I]{144}
\begin{proposition}[6.2.2]
\label{I.6.2.2-ko}
준스킴 $X$에 대하여 다음 조건들은 동치이다.
\begin{enumerate}
  \item[a)] $X$는 아르틴 스킴이다.
  \item[b)] $X$는 뇌터이고 그 바탕공간은 이산이다.
  \item[c)] $X$는 뇌터이고 그 바탕공간의 모든 점은 닫힌점이다
    ($\mathrm{T}_1$ 조건).
\end{enumerate}

이 조건들이 성립하면 $X$의 바탕공간은 유한하고, $X$의 환 $A$는
$X$의 점들의 (아르틴) 국소환들의 직접곱이다.
\end{proposition}

a)가 마지막 주장을 함의함은 알려져 있다(\cite{I-13}, p.~205,
th.~3). 이때 $A$의 모든 소 아이디얼은 극대 아이디얼이고, $A$의 한
국소환 성분의 극대 아이디얼의 역상이다. 따라서 공간 $X$는 유한하고
이산이다. 그러므로 a)는 b)를 함의하고, b)가 c)를 함의함은 자명하다.
c)가 a)를 함의함을 보이기 위해 먼저 $X$가 유한함을 보이자. 실제로
$X$가 아핀인 경우로 환원할 수 있고, 모든 소 아이디얼이 극대
아이디얼인 뇌터 환은 아르틴 환임이 알려져 있다(\cite{I-13}, p.~203).
따라서 주장이 따른다. 이제 $X$의 바탕공간은 이산이고, 유한개의 점
$x_i$의 위상적 합이다. 국소환 $\mathscr{O}_{x_i}=A_i$들은 아르틴
환이며, $A\simeq\prod_i A_i$이고 $X$는 이 환 $A$의 소 스펙트럼인
아핀 스킴과 동형임이 자명하다(\hyperref[I.1.7.3-ko]{1.7.3}).

\subsection{유한형 사상}
\label{subsection:I.6.3-ko}

\begin{definition}[6.3.1]
\label{I.6.3.1-ko}
$Y$가 다음 성질을 갖는 아핀 열린집합들의 족 $(V_\alpha)$의 합집합이면
사상 $f:X\to Y$를 \emph{유한형}이라고 한다.
\begin{enumerate}
  \item[(P)] $f^{-1}(V_\alpha)$는 유한개의 아핀 열린집합
    $U_{\alpha i}$의 합집합이고, 각 환 $A(U_{\alpha i})$는
    $A(V_\alpha)$ 위의 유한 생성 대수이다.
\end{enumerate}
이때 $X$를 $Y$ 위의 유한형 준스킴 또는 유한형 $Y$-준스킴이라고도
한다.
\end{definition}

\begin{proposition}[6.3.2]
\label{I.6.3.2-ko}
$f:X\to Y$가 유한형 사상이면 $Y$의 모든 아핀 열린집합 $W$는
정의 (\hyperref[I.6.3.1-ko]{6.3.1})의 성질 (P)를 갖는다.
\end{proposition}

먼저 다음 보조정리를 증명하자.

\begin{lemma}[6.3.2.1]
\label{I.6.3.2.1-ko}
$T\subset Y$가 성질 (P)를 갖는 아핀 열린집합이면, 모든
$g\in A(T)$에 대하여 $D(g)$도 성질 (P)를 갖는다.
\end{lemma}

실제로 가정에 따라 $f^{-1}(T)$는 유한개의 아핀 열린집합 $Z_j$의
합집합이고, 각 $A(Z_j)$는 $A(T)$ 위의 유한 생성 대수이다. $f$를
$Z_j$에 제한한 사상에 대응하는 환 준동형을
$\varphi_j:A(T)\to A(Z_j)$라 하고
(\hyperref[I.2.2.4-ko]{2.2.4}), $g_j=\varphi_j(g)$라 놓자. 그러면
$f^{-1}(D(g))\cap Z_j=D(g_j)$이다
(\hyperref[I.1.2.2.2-ko]{1.2.2.2}). 한편
$A(D(g_j))=A(Z_j)_{g_j}=A(Z_j)[1/g_j]$는 $A(Z_j)$ 위의 유한 생성
대수이고, 따라서 가정에 의해 \emph{특히} $A(T)$ 위의 유한 생성
대수이며, 그러므로 $A(D(g))=A(T)[1/g]$ 위의 유한 생성 대수이기도 하다.
이로써 보조정리가 증명된다.

이제 $W$가 준콤팩트이므로(\hyperref[I.1.1.10-ko]{1.1.10}), 각
$g_i$가 어떤 환 $A(V_{\alpha(i)})$에 속하는
$D(g_i)\subset W$들로 이루어진 $W$의 유한 덮개가 존재한다. 각
$D(g_i)$도 준콤팩트이므로
$h_{ik}\in A(W)$인 집합 $D(h_{ik})$들의 유한 합집합이다. 표준 사상을
$\varphi_i:A(W)\to A(D(g_i))$라 하면
(\hyperref[I.1.2.2.2-ko]{1.2.2.2})에 따라
$D(h_{ik})=D(\varphi_i(h_{ik}))$이다.
(\hyperref[I.6.3.2.1-ko]{6.3.2.1})에 따라 각
$f^{-1}(D(h_{ik}))$에는 유한 아핀 열린 덮개 $(U_{ijk})$가 존재하고,
각 $A(U_{ijk})$는
$A(D(h_{ik}))=A(W)[1/h_{ik}]$ 위의 유한 생성 대수이다. 이로부터 명제가
따른다.

\oldpage[I]{145}
따라서 $Y$ 위의 유한형 준스킴이라는 개념은 \emph{$Y$ 위에서
국소적}이라고 할 수 있다.

\begin{proposition}[6.3.3]
\label{I.6.3.3-ko}
$X$, $Y$를 두 아핀 스킴이라 하자. $X$가 $Y$ 위에서 유한형이기 위한
필요충분조건은 $A(X)$가 $A(Y)$ 위의 유한 생성 대수인 것이다.
\end{proposition}

이 조건이 충분함은 자명하므로 필요함을 증명하자. $A=A(Y)$,
$B=A(X)$라 놓자. (\hyperref[I.6.3.2-ko]{6.3.2})에 따라 $X$에는 유한
아핀 열린 덮개 $(V_i)$가 존재하여, 각 환 $A(V_i)$는 $A$ 위의 유한 생성
대수이다. 또한 $V_i$들이 준콤팩트이므로 각각은 유한개의 열린집합
$D(g_{ij})\subset V_i$로 덮이며, 여기서 $g_{ij}\in B$이다. 표준 포함
사상 $V_i\to X$에 대응하는 준동형을
$\varphi_i:B\to A(V_i)$라 하면
\[
  B_{g_{ij}}=(A(V_i))_{\varphi_i(g_{ij})}
  =A(V_i)[1/\varphi_i(g_{ij})].
\]
따라서 $B_{g_{ij}}$는 $A$ 위의 유한 생성 대수이다. 그러므로
$V_i=D(g_i)$이고 $g_i\in B$인 경우로 환원할 수 있다. 가정에 따라
$B_{g_i}$가 $A$ 위에서, $b_i$가 유한집합 $F_i\subset B$를 움직일 때의
원소 $b_i/g_i^{n_i}$들로 생성되도록 하는 정수 $n_i\geq0$가 존재한다.
$g_i$들이 유한개뿐이므로 모든 $n_i$가 같은 정수 $n$이라고 해도 된다.
또한 $D(g_i)$들이 $X$를 덮으므로 $g_i$들이 $B$에서 생성하는
아이디얼은 $B$이다. 다시 말해
$\sum_i h_i g_i=1$인 $h_i\in B$들이 존재한다. 이제 $F_i$들, $g_i$들,
$h_i$들의 합집합인 $B$의 유한 부분집합을 $F$라 하자. $B$의 부분환
$B'=A[F]$가 $B$와 같음을 보이겠다. 가정에 따라 모든 $b\in B$와 모든
$i$에 대하여 $B_{g_i}$에서 $b$의 표준상은 $b'_i/g_i^{m_i}$ 꼴이고,
여기서 $b'_i\in B'$이다. $b'_i$에 $g_i$의 적당한 거듭제곱을 곱하여
모든 $m_i$가 같은 정수 $m$이라고 해도 된다. 따라서 분수환의 정의에
의해 $N\geq m$이고 모든 $i$에 대하여 $g_i^N b\in B'$인 정수
$N$이 존재한다($N$은 $b$에 의존한다). 한편 \emph{환 $B'$ 안에서}
$g_i^N$들은 아이디얼 $B'$를 생성한다. 실제로 $g_i$들이 그러하고
$h_i\in B'$이기 때문이다. 따라서
$\sum_i c_i g_i^N=1$인 $c_i\in B'$들이 존재하고,
$b=\sum_i c_i g_i^N b\in B'$이다. 이로써 증명이 끝난다.

\begin{proposition}[6.3.4]
\label{I.6.3.4-ko}
\begin{enumerate}
  \item[(i)] 모든 닫힌 몰입 사상은 유한형이다.
  \item[(ii)] 두 유한형 사상의 합성은 유한형이다.
  \item[(iii)] $f:X\to X'$, $g:Y\to Y'$가 두 유한형 $S$-사상이면
    $f\times_S g:X\times_S Y\to X'\times_S Y'$도 유한형이다.
  \item[(iv)] $f:X\to Y$가 유한형 $S$-사상이면, 밑준스킴의 모든 확장
    $g:S'\to S$에 대하여 $f_{(S')}:X_{(S')}\to Y_{(S')}$도 유한형이다.
  \item[(v)] 두 사상의 합성 $g\circ f$가 유한형이고 $g$가 분리
    사상이면, $f$는 유한형이다.
  \item[(vi)] 사상 $f$가 유한형이면 $f_{\mathrm{red}}$도 유한형이다.
\end{enumerate}
\end{proposition}

(\hyperref[I.5.5.12-ko]{5.5.12})에 따라 (i), (ii), (iv)만 증명하면
충분하다.

(i)을 보이기 위해서는 $X$가 $Y$의 닫힌 부분준스킴이고
$X\to Y$가 표준 포함 사상인 경우만 다루면 된다. 또한
(\hyperref[I.6.3.2-ko]{6.3.2})에 따라 $Y$가 아핀이라고 해도 된다.
이때 $X$도 아핀이고(\hyperref[I.4.2.3-ko]{4.2.3}), 그 환은
$A/\mathfrak{J}$ 꼴의 몫환과 동형이다. 여기서 $A$는 $Y$의 환이고
$\mathfrak{J}$는 $A$의 아이디얼이다. $A/\mathfrak{J}$는 $A$ 위의
유한 생성 대수이므로 결론이 따른다.

이제 (ii)를 증명하자. $f:X\to Y$, $g:Y\to Z$를 두 유한형 사상이라
하고, $U$를 $Z$의 아핀 열린집합이라 하자. $g^{-1}(U)$에는 유한 아핀
열린 덮개 $(V_i)$가 존재하여 각 $A(V_i)$는 $A(U)$ 위의 유한 생성
대수이다(\hyperref[I.6.3.2-ko]{6.3.2}). 마찬가지로

\oldpage[I]{146}
각 $f^{-1}(V_i)$에는 유한 아핀 열린 덮개 $(W_{ij})$가 존재하여
$A(W_{ij})$는 $A(V_i)$ 위의 유한 생성 대수이고, 따라서 $A(U)$ 위의
유한 생성 대수이기도 하다. 이로부터 결론이 따른다.

마지막으로 (iv)를 증명하기 위해 $S=Y$인 경우만 다루면 된다. 실제로
$f$를 $Y$-사상으로 보고 밑준스킴의 확장을 $Y_{(S')}\to Y$로 보면
$f_{(S')}$는 $f_{(Y_{(S')})}$와도 같다
(\hyperref[I.3.3.9-ko]{3.3.9}). 이제 $p$, $q$를 각각 사영
$X_{(S')}\to X$, $X_{(S')}\to S'$라 하자. $V$를 $S$의 아핀
열린집합이라 하자. $f^{-1}(V)$는 유한개의 아핀 열린집합 $W_i$의
합집합이고, 각 $A(W_i)$는 $A(V)$ 위의 유한 생성 대수이다
(\hyperref[I.6.3.2-ko]{6.3.2}). $V'$를 $g^{-1}(V)$에 포함되는 $S'$의
아핀 열린집합이라 하자. $f\circ p=g\circ q$이므로 $q^{-1}(V')$는
$p^{-1}(W_i)$들의 합집합에 포함된다. 한편 교집합
$p^{-1}(W_i)\cap q^{-1}(V')$는 곱 $W_i\times_V V'$와 동일시된다
(\hyperref[I.3.2.7-ko]{3.2.7}). 이것은 환
$A(W_i)\otimes_{A(V)}A(V')$와 동형인 환을 갖는 아핀 스킴이다
(\hyperref[I.3.2.2-ko]{3.2.2}). 이 환은 가정에 따라 $A(V')$ 위의
유한 생성 대수이므로 명제가 증명된다.

\begin{corollary}[6.3.5]
\label{I.6.3.5-ko}
$f:X\to Y$를 몰입 사상이라 하자. $Y$의 바탕공간이 국소 뇌터이거나
$X$의 바탕공간이 뇌터이면 $f$는 유한형이다.
\end{corollary}

(\hyperref[I.6.3.2-ko]{6.3.2})에 따라 언제나 $Y$가 아핀이라고 해도
된다. $Y$의 바탕공간이 국소 뇌터이면 더 나아가 뇌터라고 해도 되고,
그 부분공간인 $X$의 바탕공간도 뇌터이다. 따라서 $Y$가 아핀이고
$X$의 바탕공간이 뇌터인 경우만 다루면 된다. 이때 $X$에는
$g_i\in A(Y)$인 유한개의 아핀 열린집합 $D(g_i)\subset Y$로 된 덮개가
존재하여, 각 $X\cap D(g_i)$는 $D(g_i)$에서 닫혀 있고 따라서 아핀
스킴이다(\hyperref[I.4.2.3-ko]{4.2.3}). 그런 덮개가 존재하는 것은 $X$가
$Y$에서 국소 닫혀 있기 때문이다(\hyperref[I.4.1.3-ko]{4.1.3}).
(\hyperref[I.6.3.4-ko]{6.3.4, (i)})와
(\hyperref[I.6.3.3-ko]{6.3.3})에 따라 $A(X\cap D(g_i))$는
$A(D(g_i))$ 위의 유한 생성 대수이다. 또한
$A(D(g_i))=A(Y)_{g_i}=A(Y)[1/g_i]$는 $A(Y)$ 위의 유한 생성 대수이므로
증명이 끝난다.

\begin{corollary}[6.3.6]
\label{I.6.3.6-ko}
$f:X\to Y$, $g:Y\to Z$를 두 사상이라 하자. $g\circ f$가 유한형이고,
$X$가 뇌터이거나 $X\times_Z Y$가 국소 뇌터이면 $f$는 유한형이다.
\end{corollary}

이는 (\hyperref[I.5.5.12-ko]{5.5.12})의 증명과 몰입 사상
$\Gamma_f$에 적용한 (\hyperref[I.6.3.5-ko]{6.3.5})에서 곧바로 따른다.

\begin{proposition}[6.3.7]
\label{I.6.3.7-ko}
$f:X\to Y$를 유한형 사상이라 하자. $Y$가 뇌터이면 $X$도 뇌터이고,
$Y$가 국소 뇌터이면 $X$도 국소 뇌터이다.
\end{proposition}

$Y$가 뇌터인 경우만 증명하면 된다. 이때 $Y$는 환 $A(V_i)$가 뇌터
환인 유한개의 아핀 열린집합 $V_i$의 합집합이다.
(\hyperref[I.6.3.2-ko]{6.3.2})에 따라 각 $f^{-1}(V_i)$는 유한개의
아핀 열린집합 $W_{ij}$의 합집합이고, 각 $A(W_{ij})$는 $A(V_i)$ 위의
유한 생성 대수이므로 뇌터 환이다. 따라서 $X$는 뇌터이다.

\begin{corollary}[6.3.8]
\label{I.6.3.8-ko}
$X$를 $S$ 위의 유한형 준스킴이라 하자. $S'$가 뇌터인(각각 국소
뇌터인) 밑준스킴의 모든 확장 $S'\to S$에 대하여 $X_{(S')}$는
뇌터이다(각각 국소 뇌터이다).
\end{corollary}

이는 (\hyperref[I.6.3.7-ko]{6.3.7})에서 따른다. 실제로
(\hyperref[I.6.3.4-ko]{6.3.4, (iv)})에 따라 $X_{(S')}$는 $S'$ 위에서
유한형이다.

또한 $S$-준스킴의 곱 $X\times_S Y$에서, 인자 $X$, $Y$ 가운데
\emph{하나가}
\oldpage[I]{147}
$S$ 위에서 \emph{유한형}이고
\emph{다른 하나가 뇌터이면}(각각 \emph{국소 뇌터이면}),
$X\times_S Y$는 \emph{뇌터이다}(각각 \emph{국소 뇌터이다}).

\begin{corollary}[6.3.9]
\label{I.6.3.9-ko}
$X$를 국소 뇌터 준스킴 $S$ 위의 유한형 준스킴이라 하자. 그러면 모든
$S$-사상 $f:X\to Y$는 유한형이다.
\end{corollary}

실제로 $S$가 뇌터라고 해도 된다. $\varphi:X\to S$,
$\psi:Y\to S$가 구조 사상이면 $\varphi=\psi\circ f$이고,
(\hyperref[I.6.3.7-ko]{6.3.7})에 따라 $X$는 뇌터이다. 따라서
(\hyperref[I.6.3.6-ko]{6.3.6})에 따라 $f$는 유한형이다.

\begin{proposition}[6.3.10]
\label{I.6.3.10-ko}
$f:X\to Y$를 유한형 사상이라 하자. $f$가 전사이기 위한 필요충분조건은
모든 대수적으로 닫힌 체 $\Omega$에 대하여 $f$에 대응하는 사상
$X(\Omega)\to Y(\Omega)$가 전사인 것이다
(\hyperref[I.3.4.1-ko]{3.4.1}).
\end{proposition}

이 조건이 충분함은 각 $y\in Y$에 대하여 $k(y)$의 대수적으로 닫힌
확대체 $\Omega$와 다음 가환 그림을 고려하면 알 수 있다.
\[
  \xymatrix{
    X\ar[dd]_f\\
    & \operatorname{Spec}(\Omega)\ar[ul]\ar[dl]\\
    Y
  }
\]
(\emph{참조} (\hyperref[I.3.5.3-ko]{3.5.3})). 역으로 $f$가 전사라고
가정하고, $\Omega$가 대수적으로 닫힌 체인 사상
$g:\{\xi\}=\operatorname{Spec}(\Omega)\to Y$를 택하자. 다음 그림을
고려하면
\[
  \xymatrix{
    X\ar[d]_f &
    X_{(\Omega)}\ar[l]\ar[d]^{f_{(\Omega)}}\\
    Y &
    \operatorname{Spec}(\Omega)\ar[l]
  }
\]
$X_{(\Omega)}$에 \emph{$\Omega$-유리점}이 존재함을 보이면 충분하다
(\hyperref[I.3.3.14-ko]{3.3.14}, \hyperref[I.3.4.3-ko]{3.4.3},
\hyperref[I.3.4.4-ko]{3.4.4}). $f$가 전사이므로 $X_{(\Omega)}$는
공집합이 아니고(\hyperref[I.3.5.10-ko]{3.5.10}), $f$가 유한형이므로
(\hyperref[I.6.3.4-ko]{6.3.4, (iv)})에 따라 $f_{(\Omega)}$도 유한형이다.
따라서 $X_{(\Omega)}$에는 $A(Z)$가 $\Omega$ 위의 영환이 아닌 유한 생성
대수인 공집합이 아닌 아핀 열린집합 $Z$가 존재한다. 힐베르트
영점정리(\cite{I-21})에 따라 $\Omega$-대수 준동형
$A(Z)\to\Omega$가 존재하고, 따라서 $\operatorname{Spec}(\Omega)$ 위의
$X_{(\Omega)}$의 단면이 존재한다. 이로써 명제가 증명된다.

\subsection{대수적 준스킴}
\label{subsection:I.6.4-ko}

\begin{definition}[6.4.1]
\label{I.6.4.1-ko}
체 $K$가 주어졌을 때 $K$ 위의 유한형 준스킴 $X$를
\emph{$K$-대수적 준스킴}이라고 한다. $K$를 $X$의 \emph{기초체}라고 한다.
더 나아가 $X$가 스킴이면(또는 이와 동치로
(\hyperref[I.5.5.8-ko]{5.5.8}), $X$가 $K$-스킴이면) $X$를
\emph{$K$-대수적 스킴}이라고도 한다.
\end{definition}

모든 $K$-대수적 준스킴은 \emph{뇌터이다}
(\hyperref[I.6.3.7-ko]{6.3.7}).

\begin{proposition}[6.4.2]
\label{I.6.4.2-ko}
$X$를 $K$-대수적 준스킴이라 하자. 점 $x\in X$가 닫힌점이기 위한
필요충분조건은 $k(x)$가 $K$의 유한 차수 대수적 확대체인 것이다.
\end{proposition}

$X$가 아핀이고 그 환 $A$가 $K$ 위의 유한 생성 대수인 경우만 다루면
된다. 실제로 $A(U)$가 $K$ 위의 유한 생성 대수인 $X$의 아핀 열린집합
$U$들은 $X$를 덮는다(\hyperref[I.6.3.1-ko]{6.3.1}). 이때 $X$의 닫힌점들은
$\mathfrak{j}_x$가
\oldpage[I]{148}
$A$의 극대 아이디얼인 점들, 다시 말해 $A/\mathfrak{j}_x$가 체인 점들이다.
이 체는 반드시 $k(x)$와 같다. $A/\mathfrak{j}_x$가 $K$ 위의 유한 생성
대수이므로, $x$가 닫힌점이면 $k(x)$는 $K$ 위의 유한 생성 대수인 체이고,
따라서 반드시 $K$ 위에서 \emph{유한 차원}이다 \cite{I-21}. 역으로
$k(x)$가 $K$ 위에서 유한 차원이면 그 부분환
$A/\mathfrak{j}_x\subset k(x)$도 그러하다. $K$ 위의 유한 차원 대수인
정역은 모두 체이므로 $A/\mathfrak{j}_x=k(x)$이고, 따라서 $x$는 닫힌점이다.

\begin{corollary}[6.4.3]
\label{I.6.4.3-ko}
$K$를 대수적으로 닫힌 체, $X$를 $K$-대수적 준스킴이라 하자. 그러면
$X$의 닫힌점들은 $K$-유리점들이고(\hyperref[I.3.4.4-ko]{3.4.4}),
$K$에 값을 갖는 $X$의 점들과 표준적으로 동일시된다.
\end{corollary}

\begin{proposition}[6.4.4]
\label{I.6.4.4-ko}
$X$를 체 $K$ 위의 대수적 준스킴이라 하자. 다음 성질들은 동치이다.
\begin{enumerate}
  \item[\rm(a)] $X$는 아르틴이다.
  \item[\rm(b)] $X$의 바탕공간은 이산공간이다.
  \item[\rm(c)] $X$의 바탕공간에는 닫힌점이 유한개뿐이다.
  \item[\rm(c')] $X$의 바탕공간은 유한집합이다.
  \item[\rm(d)] $X$의 모든 점은 닫힌점이다.
  \item[\rm(e)] $X$는 $K$ 위의 유한 차원 대수 $A$에 대하여
    $\operatorname{Spec}(A)$와 동형이다.
\end{enumerate}
\end{proposition}

$X$가 뇌터이므로 (\hyperref[I.6.2.2-ko]{6.2.2})에 따라 조건 (a), (b),
(d)는 동치이고 (c)와 (c')를 함의한다. 또한 (e)가 (a)를 함의함은
명백하다. (c)가 (d)와 (e)를 함의함을 보이는 일만 남는다. $X$가 아핀인
경우만 다루면 된다. 이때 $A(X)$는 $K$ 위의 유한 생성 대수이고
(\hyperref[I.6.3.3-ko]{6.3.3}), 따라서 야콥슨 환이다
(\cite[p.~3-11 및 3-12]{I-1}). 가정에 따라 이 환에는 극대 아이디얼이
유한개뿐이다. 소 아이디얼들의 유한 교집합이 소 아이디얼이면 그 가운데
하나와 같아야 하므로, $A(X)$의 모든 소 아이디얼은 극대 아이디얼이다.
따라서 (d)가 따른다. 또한 (\hyperref[I.6.2.2-ko]{6.2.2})에 따라
$A(X)$는 유한 생성 아르틴 $K$-대수이고, 따라서 반드시
\emph{유한 차원}이다 \cite{I-21}.

\begin{env}[6.4.5]
\label{I.6.4.5-ko}
(\hyperref[I.6.4.4-ko]{6.4.4})의 조건들이 성립할 때 $X$를
\emph{$K$ 위의 유한 스킴}(\emph{참조}
(\agref{II.6.1.1-ko}{II, 6.1.1})) 또는 \emph{유한 $K$-스킴}이라고 한다.
그 \emph{계수(階數)}는 $[A:K]$이며 $\operatorname{rg}_K(X)$로도 쓴다.
$X$, $Y$가 $K$ 위의 두 유한 스킴이면
\[
  \operatorname{rg}_K(X\amalg Y)
  =\operatorname{rg}_K(X)+\operatorname{rg}_K(Y)
  \tag{6.4.5.1}\label{I.6.4.5.1-ko}
\]
\[
  \operatorname{rg}_K(X\times_K Y)
  =\operatorname{rg}_K(X)\operatorname{rg}_K(Y)
  \tag{6.4.5.2}\label{I.6.4.5.2-ko}
\]
이다. 이는 (\hyperref[I.3.2.2-ko]{3.2.2})에서 따른다.
\end{env}

\begin{corollary}[6.4.6]
\label{I.6.4.6-ko}
$X$를 체 $K$ 위의 유한 스킴이라 하자. $K$의 모든 확대체 $K'$에 대하여
$X\otimes_K K'$는 $K'$ 위의 유한 스킴이고, $K'$ 위의 계수는 $X$의
$K$ 위의 계수와 같다.
\end{corollary}

실제로 $A=A(X)$라 하면 $[A\otimes_K K':K']=[A:K]$이다.

\begin{corollary}[6.4.7]
\label{I.6.4.7-ko}
$X$를 체 $K$ 위의 유한 스킴이라 하고
$n=\sum_{x\in X}[k(x):K]_s$라 놓자.
\emph{(확대체 $K'/K$에 대하여 $[K':K]_s$는 $K'$에 포함되는 $K$의
가장 큰 대수적 분리확대체의 $K$ 위의 \emph{분리차수}임을 상기하자.)}
\oldpage[I]{149}
그러면 $K$의 모든 대수적으로 닫힌 확대체 $\Omega$에 대하여
$X\otimes_K\Omega$의 바탕공간에는 점이 정확히 $n$개 있고, 이 점들은
$X(\Omega)_K$의 원소들과 동일시된다.
\end{corollary}

(\hyperref[I.6.2.2-ko]{6.2.2})에 따라 환 $A=A(X)$가 국소환인 경우만
다루면 된다. $\mathfrak m$을 그 극대 아이디얼, $L=A/\mathfrak m$을
$K$의 대수적 확대체인 잉여체라 하자. 그러면 $X(\Omega)_K$의 원소들은
$X\otimes_K\Omega$의 $\Omega$-단면들과 전단사로 대응하고
(\hyperref[I.3.4.1-ko]{3.4.1}, \hyperref[I.3.3.14-ko]{3.3.14}), 또한
$L$에서 $\Omega$로 가는 $K$-준동형들과 전단사로 대응한다
(\hyperref[I.1.7.3-ko]{1.7.3}). 따라서
(Bourbaki, \emph{Alg.}, 제V장, \textsection7, n\textsuperscript{o}~5,
명제~8)과 (\hyperref[I.6.4.3-ko]{6.4.3})을 고려하면 결론이 따른다.

\begin{env}[6.4.8]
\label{I.6.4.8-ko}
(\hyperref[I.6.4.7-ko]{6.4.7})에서 정의한 수 $n$을 $A$(또는 $X$)의
$K$ 위의 \emph{분리차수}, 또는 \emph{$X$의 점의 기하적 개수}라고 한다.
따라서 이는 $X(\Omega)_K$의 원소 수와 같다. 이 정의에서 곧바로,
$K$의 모든 확대체 $K'$에 대하여 $X\otimes_K K'$의 점의 기하적 개수가
$X$의 점의 기하적 개수와 같음이 따른다. 이 수를 $n(X)$라 하자. $X$, $Y$가
$K$ 위의 두 유한 스킴이면 명백히
\[
  n(X\amalg Y)=n(X)+n(Y).
  \tag{6.4.8.1}\label{I.6.4.8.1-ko}
\]
이다.

같은 가정 아래에서
\[
  n(X\times_K Y)=n(X)n(Y)
  \tag{6.4.8.2}\label{I.6.4.8.2-ko}
\]
도 성립한다. 이는 $n(X)$를 $X(\Omega)_K$의 원소 수로 해석하고 공식
(\hyperref[I.3.4.3.1-ko]{3.4.3.1})을 적용하면 곧바로 따른다.
\end{env}

\begin{proposition}[6.4.9]
\label{I.6.4.9-ko}
$K$를 체, $X$, $Y$를 두 $K$-대수적 준스킴,
$f:X\to Y$를 $K$-사상이라 하자. 또한 $\Omega$를 $K$ 위의 초월차수가
무한인 대수적으로 닫힌 확대체라 하자. $f$가 전사이기 위한 필요충분조건은
$f$에 대응하는 사상 $X(\Omega)_K\to Y(\Omega)_K$가 전사인 것이다
(\hyperref[I.3.4.1-ko]{3.4.1}).
\end{proposition}

필요성은 (\hyperref[I.6.3.10-ko]{6.3.10})에서 따른다. 여기서 $f$가
반드시 유한형인 것은 (\hyperref[I.6.3.9-ko]{6.3.9})에 따른다.
충분성을 보이기 위해 (\hyperref[I.6.3.10-ko]{6.3.10})에서와 같이
논증한다. 실제로 모든 $y\in Y$에 대하여 $k(y)$는 $K$의 유한 생성
확대체이고, 따라서 $\Omega$의 부분체와 $K$-동형이다.

\begin{remark}[6.4.10]
\label{I.6.4.10-ko}
제IV장에서 (\hyperref[I.6.4.9-ko]{6.4.9})의 결론이 $\Omega$의 $K$ 위의
초월차수에 관한 가정 없이도 성립함을 보일 것이다.
\end{remark}

\begin{proposition}[6.4.11]
\label{I.6.4.11-ko}
$f:X\to Y$가 유한형 사상이면, 모든 $y\in Y$에 대하여 섬유
$f^{-1}(y)$는 잉여체 $k(y)$ 위의 대수적 준스킴이고, 모든
$x\in f^{-1}(y)$에 대하여 $k(x)$는 $k(y)$의 유한 생성 확대체이다.
\end{proposition}

$f^{-1}(y)=X\otimes_Y k(y)$이므로(\hyperref[I.3.6.3-ko]{3.6.3}), 이 명제는
(\hyperref[I.6.3.4-ko]{6.3.4, (iv)})와
(\hyperref[I.6.3.3-ko]{6.3.3})에서 따른다.

\begin{proposition}[6.4.12]
\label{I.6.4.12-ko}
$f:X\to Y$, $g:Y'\to Y$를 두 사상이라 하자.
$X'=X\times_Y Y'$라 놓고 $f'=f_{(Y')}:X'\to Y'$라 하자.
$y'\in Y'$, $y=g(y')$라 하자. 섬유 $f^{-1}(y)$가 $k(y)$ 위의 유한
대수적 스킴이면, 섬유 $f'^{-1}(y')$도 $k(y')$ 위의 유한 대수적
스킴이고, 그 계수와 점의 기하적 개수는 $f^{-1}(y)$의 것과 같다.
\end{proposition}

섬유의 추이성(\hyperref[I.3.6.5-ko]{3.6.5})을 고려하면 이는
(\hyperref[I.6.4.6-ko]{6.4.6})과 (\hyperref[I.6.4.8-ko]{6.4.8})에서
곧바로 따른다.

\begin{env}[6.4.13]
\label{I.6.4.13-ko}
명제 (\hyperref[I.6.4.11-ko]{6.4.11})은 유한형 사상이 직관적으로
``대수다양체들의 대수적 족''에 해당함을 보여 준다. 여기서 $Y$의 점들은
\oldpage[I]{150}
``매개변수''의 역할을 하며, 이것이 이 사상들에 ``기하적'' 의미를 준다.
유한형이 아닌 사상들은 뒤에서 주로, 예를 들어 국소화나 완비화에 의한
``밑준스킴의 변경'' 문제에 나타날 것이다.
\end{env}

\subsection{사상의 국소적 결정}
\label{subsection:I.6.5-ko}

\begin{proposition}[6.5.1]
\label{I.6.5.1-ko}
$X$, $Y$를 두 $S$-준스킴이라 하고, $Y$가 $S$ 위의 유한형이라고 하자.
또한 $x\in X$, $y\in Y$가 같은 점 $s\in S$ 위에 있다고 하자.
\begin{enumerate}
  \item[\rm(i)] $X$에서 $Y$로 가는 두 $S$-사상
    $f=(\psi,\theta)$, $f'=(\psi',\theta')$가
    $\psi(x)=\psi'(x)=y$를 만족하고, $\mathscr{O}_y$에서
    $\mathscr{O}_x$로 가는 (국소) $\mathscr{O}_s$-준동형
    $\theta_x^\sharp$와 ${\theta'}_x^\sharp$가 같으면, $f$와 $f'$는
    $x$의 어떤 열린 근방에서 일치한다.
  \item[\rm(ii)] 더 나아가 $S$가 국소 뇌터라고 하자. 모든 국소
    $\mathscr{O}_s$-준동형
    $\varphi:\mathscr{O}_y\to\mathscr{O}_x$에 대하여 $X$에서 $x$의
    열린 근방 $U$와 $U$에서 $Y$로 가는 $S$-사상
    $f=(\psi,\theta)$가 존재하여
    $\psi(x)=y$이고 $\theta_x^\sharp=\varphi$이다.
\end{enumerate}
\end{proposition}

\begin{enumerate}
  \item[\rm(i)] 문제는 $S$, $X$, $Y$ 위에서 국소적이므로 $S$, $X$, $Y$가
    각각 환 $A$, $B$, $C$의 아핀 준스킴이라고 해도 된다. 이때 $f$, $f'$는
    각각 $({}^a\varphi,\widetilde{\varphi})$,
    $({}^a\varphi',\widetilde{\varphi}')$의 꼴이며, $\varphi$, $\varphi'$는
    $C$에서 $B$로 가는 두 $A$-준동형으로서
    $\varphi^{-1}(\mathfrak{j}_x)
    ={\varphi'}^{-1}(\mathfrak{j}_x)=\mathfrak{j}_y$를 만족한다. 또한
    $\varphi$, $\varphi'$에서 얻어지는 $C_y$에서 $B_x$로
    가는 준동형 $\varphi_x$, $\varphi'_x$는 같다. 더 나아가 $C$가 $A$ 위의
    유한 생성 대수라고 해도 된다. $c_i$($1\leq i\leq n$)를 $A$-대수 $C$의
    생성원이라 하고 $b_i=\varphi(c_i)$, $b'_i=\varphi'(c_i)$라 놓자. 가정에
    따라 분수환 $B_x$에서 $b_i/1=b'_i/1$이다($1\leq i\leq n$). 이는
    $1\leq i\leq n$에 대하여 $s_i(b_i-b'_i)=0$이 되는 원소
    $s_i\in B-\mathfrak{j}_x$가 존재한다는 뜻이며, 모든 $s_i$가 같은 원소
    $g\in B-\mathfrak{j}_x$라고 해도 된다. 따라서 분수환 $B_g$에서
    $1\leq i\leq n$에 대하여 $b_i/1=b'_i/1$이다. $i_g:B\to B_g$를 표준
    준동형이라 하면 $i_g\circ\varphi=i_g\circ\varphi'$이고, 그러므로
    $f$와 $f'$의 $D(g)$로의 제한은 같다.
  \item[\rm(ii)] (i)와 같은 상황으로 환원할 수 있고, 더 나아가 환 $A$가
    뇌터라고 해도 된다. $c_i$($1\leq i\leq n$)를 $A$-대수 $C$의 생성원이라
    하고, $\alpha:A[X_1,\ldots,X_n]\to C$를 $X_i$를 $c_i$로 보내는
    다항식 대수 $A[X_1,\ldots,X_n]$에서 $C$로 가는 전사 준동형이라 하자
    ($1\leq i\leq n$). 한편 $i_y:C\to C_y$를 표준 준동형이라 하고, 합성
    준동형
    \[
      \beta:A[X_1,\ldots,X_n]
      \xrightarrow{\alpha} C
      \xrightarrow{i_y} C_y
      \xrightarrow{\varphi} B_x
    \]
    을 생각하자. $\mathfrak{a}$를 $\beta$의 핵이라 하자. $A$가 뇌터이므로
    $A[X_1,\ldots,X_n]$도 뇌터이고, 따라서 $\mathfrak{a}$는 유한 생성계
    $Q_j(X_1,\ldots,X_n)$($1\leq j\leq m$)를 갖는다. 한편 각 원소
    $\varphi(i_y(c_i))$는 $b_i\in B$, $s_i\notin\mathfrak{j}_x$에 대하여
    $b_i/s_i$로 쓸 수 있고, 모든 $s_i$가 같은 원소
    $g\in B-\mathfrak{j}_x$라고 해도 된다. 이제 가정에 따라 $B_x$에서
    $Q_j(b_1/g,\ldots,b_n/g)=0$이다. 다음과 같이 놓자.
    \[
      Q_j(X_1/T,\ldots,X_n/T)
      =P_j(X_1,\ldots,X_n,T)/T^{k_j},
    \]
    여기서 $P_j$는 $k_j$차 동차다항식이다. 또한
    $d_j=P_j(b_1,\ldots,b_n,g)\in B$라 하자. 가정에 따라
    $t_j\in B-\mathfrak{j}_x$에 대하여 $t_jd_j=0$이다
    ($1\leq j\leq m$). 모든 $t_j$가
\oldpage[I]{151}
    같은 원소 $h\in B-\mathfrak{j}_x$라고 해도 된다. 따라서
    $1\leq j\leq m$에 대하여
    $P_j(hb_1,\ldots,hb_n,hg)=0$이다. 이제 $X_i$를 $hb_i/hg$로 보내는
    $A[X_1,\ldots,X_n]$에서 분수환 $B_{hg}$로 가는 준동형 $\rho$를
    생각하자($1\leq i\leq n$). 이 준동형에 의한 $\mathfrak a$의 상은
    0이고, 따라서 $\rho$에 의한 $\alpha^{-1}(0)$의 상도
    \emph{특히} 0이다. 그러므로 $\rho$는
    $A[X_1,\ldots,X_n]\xrightarrow{\alpha}C
    \xrightarrow{\gamma}B_{hg}$로 인수분해되며,
    $\gamma(c_i)=hb_i/hg$이다. $i_x:B_{hg}\to B_x$를 표준
    준동형이라 하면 그림
    \[
      \xymatrix{
        C\ar[r]^\gamma\ar[d]_{i_y} & B_{hg}\ar[d]^{i_x}\\
        C_y\ar[r]_\varphi & B_x
      }
      \tag{6.5.1.1}\label{I.6.5.1.1-ko}
    \]
    은 가환임이 명백하다. 따라서 $\varphi=\gamma_x$이고, $\varphi$가 국소
    준동형이므로 ${}^a\gamma(x)=y$이다. 그러므로
    $f=({}^a\gamma,\widetilde{\gamma})$는 $x$의 근방 $D(hg)$에서 $Y$로
    가는, 문제의 조건을 만족하는 $S$-사상이다.
\end{enumerate}

\begin{corollary}[6.5.2]
\label{I.6.5.2-ko}
(\hyperref[I.6.5.1-ko]{6.5.1, (ii)})의 가정 아래에서, 더 나아가 $X$가
$S$ 위의 유한형이면 사상 $f$가 유한형이라고 할 수 있다.
\end{corollary}

이는 (\hyperref[I.6.3.6-ko]{6.3.6})에서 따른다.

\begin{corollary}[6.5.3]
\label{I.6.5.3-ko}
(\hyperref[I.6.5.1-ko]{6.5.1, (ii)})의 가정들이 성립한다고 하자. 더 나아가
$Y$가 정역 준스킴이고 $\varphi$가 단사 준동형이면,
$f=({}^a\gamma,\widetilde{\gamma})$이고 $\gamma$가 단사 준동형이라고 할 수
있다.
\end{corollary}

실제로 $C$가 정역이라고 해도 되고(\hyperref[I.5.1.4-ko]{5.1.4}), 따라서
$i_y$는 단사 준동형이다. 그러므로 그림
(\hyperref[I.6.5.1.1-ko]{6.5.1.1})에서 $\gamma$가 단사 준동형임이 따른다.

\begin{proposition}[6.5.4]
\label{I.6.5.4-ko}
$f=(\psi,\theta):X\to Y$를 유한형 사상, $x$를 $X$의 점,
$y=\psi(x)$라 하자.
\begin{enumerate}
  \item[\rm(i)] $f$가 점 $x$에서 국소 몰입 사상이기 위한
    (\hyperref[I.4.5.1-ko]{4.5.1}) 필요충분조건은
    $\theta_x^\sharp:\mathscr{O}_y\to\mathscr{O}_x$가 전사 준동형인 것이다.
  \item[\rm(ii)] 더 나아가 $Y$가 국소 뇌터라고 하자. $f$가 점 $x$에서
    국소 동형사상이기 위한(\hyperref[I.4.5.2-ko]{4.5.2}) 필요충분조건은
    $\theta_x^\sharp$가 동형인 것이다.
\end{enumerate}
\end{proposition}

\begin{enumerate}
  \item[\rm(ii)] (\hyperref[I.6.5.1-ko]{6.5.1})에 따라 $y$의 열린 근방
    $V$와 사상 $g:V\to X$가 존재하여 $g\circ f$(각각 $f\circ g$)가
    정의되고 $x$(각각 $y$)의 어떤 근방에서 항등사상과 일치한다. 이로부터
    $f$가 국소 동형사상임이 곧바로 따른다.
  \item[\rm(i)] 문제는 $X$, $Y$ 위에서 국소적이므로 $X$, $Y$가 각각 환
    $A$, $B$의 아핀 준스킴이라고 해도 된다. 이때
    $f=({}^a\varphi,\widetilde{\varphi})$이고, $\varphi:B\to A$는 $A$를
    유한 생성 $B$-대수로 만드는 환 준동형이다. 또한
    $\varphi^{-1}(\mathfrak{j}_x)=\mathfrak{j}_y$이고, $\varphi$에서
    얻어지는 준동형 $\varphi_x:B_y\to A_x$는 전사이다. $(t_i)$
    ($1\leq i\leq n$)를 $B$-대수 $A$의 생성계라 하자. $\varphi_x$에 관한
    가정에 따라 $b_i\in B$와 $c\in B-\mathfrak{j}_y$가 존재하여 분수환
    $A_x$에서 $1\leq i\leq n$에 대하여
    $t_i/1=\varphi(b_i)/\varphi(c)$이다. 따라서
    (\hyperref[I.1.3.3-ko]{1.3.3})에 따라 $a\in A-\mathfrak{j}_x$가
    존재하여, $g=a\varphi(c)$라 놓으면
    $t_i/1=a\varphi(b_i)/g$가 \emph{분수환 $A_g$ 안에서}도 성립한다.
    한편 가정에 따라 $\varphi(B)$를 계수환으로 하는 다항식
    $Q(X_1,\ldots,X_n)$가
\oldpage[I]{152}
    존재하여 $a=Q(t_1,\ldots,t_n)$이다. 다음과 같이 놓자.
    \[
      Q(X_1/T,\ldots,X_n/T)
      =P(X_1,\ldots,X_n,T)/T^m,
    \]
    여기서 $P$는 $m$차 동차다항식이다. 환 $A_g$에서
    \[
      a/1
      =a^mP(\varphi(b_1),\ldots,\varphi(b_n),\varphi(c))/g^m
      =a^m\varphi(d)/g^m
    \]
    이며, 여기서 $d\in B$이다. $A_g$에서
    $g/1=(a/1)(\varphi(c)/1)$은 정의에 따라 가역이므로 $a/1$과
    $\varphi(c)/1$도 가역이고, 따라서
    \[
      a/1=(\varphi(d)/1)(\varphi(c)/1)^{-m}
    \]
    으로 쓸 수 있다. 그러므로 $\varphi(d)/1$도 $A_g$에서 가역이다.
    이제 $h=cd$라 놓자. $\varphi(h)/1$이 $A_g$에서 가역이므로 합성
    준동형
    \[
      B\xrightarrow{\varphi}A\longrightarrow A_g
    \]
    은 다음과 같이 인수분해된다.
    \[
      B\longrightarrow B_h\xrightarrow{\gamma}A_g
    \]
    (\hyperref[0.1.2.4-ko]{0, 1.2.4}). $\gamma$가
    \emph{전사}임을 보이자. $A_g$에서 $B_h$의 상이 $t_i/1$과
    $(g/1)^{-1}$을 포함함을 보이면 충분하다. 그런데
    \[
      (g/1)^{-1}
      =(\varphi(c)/1)^{m-1}(\varphi(d)/1)^{-1}
      =\gamma(c^m/h),
    \]
    이고
    \[
      a/1=\gamma(d^{m+1}/h^m),
    \]
    이므로
    \[
      (a\varphi(b_i))/1=\gamma(b_i d^{m+1}/h^m)
    \]
    이다. 또한
    $t_i/1=(a\varphi(b_i)/1)(g/1)^{-1}$이므로 주장이 증명된다. $h$의
    선택에 따라 $\psi(D(g))\subset D(h)$이고, $f$를 $D(g)$에 제한한
    사상은 $({}^a\gamma,\widetilde{\gamma})$와 같다. $\gamma$가 전사이므로
    이 제한은 $D(g)$에서 $D(h)$로 가는 닫힌 몰입 사상이다
    (\hyperref[I.4.2.3-ko]{4.2.3}).
\end{enumerate}

\begin{corollary}[6.5.5]
\label{I.6.5.5-ko}
$f=(\psi,\theta):X\to Y$를 유한형 사상이라 하자. $X$가 기약이라고 하고,
그 일반점을 $x$라 하며 $y=\psi(x)$라 놓자.
\begin{enumerate}
  \item[\rm(i)] $f$가 $X$의 어떤 점에서 국소 몰입 사상이기 위한
    필요충분조건은
    $\theta_x^\sharp:\mathscr{O}_y\to\mathscr{O}_x$가 전사 준동형인 것이다.
  \item[\rm(ii)] 더 나아가 $Y$가 기약이고 국소 뇌터라고 하자. $f$가
    $X$의 어떤 점에서 국소 동형사상이기 위한 필요충분조건은 $y$가 $Y$의
    일반점인 것, 다시 말해(\hyperref[0.2.1.4-ko]{0, 2.1.4}) $f$가
    \emph{우세 사상}인 것과 $\theta_x^\sharp$가 동형인 것, 다시 말해
    $f$가 \emph{쌍유리 사상}인 것이다
    (\hyperref[I.2.2.9-ko]{2.2.9}).
\end{enumerate}
\end{corollary}

(i)는 $X$의 공집합이 아닌 모든 열린집합이 $x$를 포함한다는 사실을
고려하면 (\hyperref[I.6.5.4-ko]{6.5.4, (i)})에서 명백히 따른다. 마찬가지로
(ii)는 (\hyperref[I.6.5.4-ko]{6.5.4, (ii)})에서 따른다.

\subsection{준콤팩트 사상과 국소 유한형 사상}
\label{subsection:I.6.6-ko}

\begin{definition}[6.6.1]
\label{I.6.6.1-ko}
사상 $f:X\to Y$에 의한 $Y$의 모든 준콤팩트 열린집합의 역상이
준콤팩트이면, $f$를 \emph{준콤팩트} 사상이라 한다.
\end{definition}

$\mathfrak{B}$를 준콤팩트 열린집합들(예를 들어 아핀 열린집합들)로
이루어진 $Y$의 위상의 기저라 하자. $f$가 준콤팩트이기 위한 필요충분조건은
$\mathfrak{B}$의 모든 원소의 $f$에 의한 역상이 준콤팩트인 것, 또는 같은
말로 아핀 열린집합들의 \emph{유한} 합집합인 것이다. 실제로 $Y$의 모든
준콤팩트 열린집합은 $\mathfrak{B}$의 유한개의 원소의 합집합이다. 예를 들어
$X$가 \emph{준콤팩트}이고 $Y$가 \emph{아핀}이면 \emph{모든} 사상
$f:X\to Y$는 준콤팩트이다. 실제로 $X$는 유한개의 아핀 열린집합 $U_i$의
합집합이고, $Y$의 모든 아핀 열린집합 $V$에 대하여
$U_i\cap f^{-1}(V)$는 아핀이므로
(\hyperref[I.5.5.10-ko]{5.5.10}) 준콤팩트이다.

$f:X\to Y$가 준콤팩트 사상이면, $Y$의 모든 열린집합 $V$에 대하여
$f$를 $f^{-1}(V)$에 제한한 사상 $f^{-1}(V)\to V$가 준콤팩트임은
명백하다. 역으로 $(U_\alpha)$가 $Y$의 열린 덮개이고, 사상 $f:X\to Y$의
제한들 $f^{-1}(U_\alpha)\to U_\alpha$가 준콤팩트이면, $f$도
준콤팩트이다.

\begin{definition}[6.6.2]
\label{I.6.6.2-ko}
사상 $f:X\to Y$가 다음 조건을 만족하면 \emph{국소 유한형} 사상이라 한다.
모든 $x\in X$에 대하여 $x$의 열린 근방 $U$와
$f(U)$를 포함하는 $y=f(x)$의 열린 근방 $V$가 존재하고, $f$의
\oldpage[I]{153}
$U$로의 제한이 $U$에서 $V$로 가는 유한형 사상이다. 이때 $X$를 $Y$
위의 국소 유한형 준스킴 또는 국소 유한형 $Y$-준스킴이라고도 한다.
\end{definition}

(\hyperref[I.6.3.2-ko]{6.3.2})에 따라 $f$가 국소 유한형이면, $Y$의 모든
열린집합 $W$에 대하여 $f$를 $f^{-1}(W)$에 제한한 사상
$f^{-1}(W)\to W$도 국소 유한형임이 곧바로 따른다.

$Y$가 국소 뇌터이고 $X$가 $Y$ 위의 국소 유한형이면,
(\hyperref[I.6.3.7-ko]{6.3.7})에 따라 $X$는 국소 뇌터이다.

\begin{proposition}[6.6.3]
\label{I.6.6.3-ko}
사상 $f:X\to Y$가 유한형이기 위한 필요충분조건은 $f$가 준콤팩트이고
국소 유한형인 것이다.
\end{proposition}

조건의 필요성은 (\hyperref[I.6.3.1-ko]{6.3.1})과
(\hyperref[I.6.6.1-ko]{6.6.1}) 뒤의 주에서 곧바로 따른다. 역으로 이
조건들이 성립한다고 하고, $U$를 $Y$의 아핀 열린집합, $A$를 그 환이라
하자. 모든
$x\in f^{-1}(U)$에 대하여 가정에 따라 $x$의 근방
$V(x)\subset f^{-1}(U)$와 $y=f(x)$의 근방 $W(x)\subset U$가 존재하여,
$W(x)$는 $f(V(x))$를 포함하고 $f$를 $V(x)$에 제한한 사상
$V(x)\to W(x)$는 유한형이다. $W(x)$를 $y=f(x)$의
$D(g)$ 꼴인 근방 $W_1(x)\subset W(x)$로($g\in A$), $V(x)$를
$V(x)\cap f^{-1}(W_1(x))$로 바꾸면, $W(x)$가 $D(g)$ 꼴이라고 해도 된다.
따라서 그 환이 $A[1/g]$로 쓰이므로 $W(x)$는 $U$ 위의 유한형이고, 그러므로
$V(x)$도 $U$ 위의 유한형이다. 더 나아가 $f^{-1}(U)$는 가정에 따라
준콤팩트이므로 유한개의 열린집합 $V(x_i)$의 합집합이다. 이로써 증명이
끝난다.

\begin{proposition}[6.6.4]
\label{I.6.6.4-ko}
\begin{enumerate}
  \item[\rm(i)] 몰입 사상 $X\to Y$가 닫힌 몰입 사상이거나, $Y$의
    바탕공간이 국소 뇌터이거나, $X$의 바탕공간이 뇌터이면 이 몰입 사상은
    준콤팩트이다.
  \item[\rm(ii)] 두 준콤팩트 사상의 합성은 준콤팩트이다.
  \item[\rm(iii)] $f:X\to Y$가 준콤팩트 $S$-사상이면, 밑준스킴의 모든
    확장 $g:S'\to S$에 대하여
    $f_{(S')}:X_{(S')}\to Y_{(S')}$도 준콤팩트이다.
  \item[\rm(iv)] $f:X\to X'$와 $g:Y\to Y'$가 두 준콤팩트 $S$-사상이면
    $f\times_S g$도 준콤팩트이다.
  \item[\rm(v)] 두 사상 $f:X\to Y$, $g:Y\to Z$의 합성이 준콤팩트이고,
    또한 $g$가 분리 사상이거나 $X$의 바탕공간이 국소 뇌터이면 $f$는
    준콤팩트이다.
  \item[\rm(vi)] 사상 $f$가 준콤팩트이기 위한 필요충분조건은
    $f_{\mathrm{red}}$가 준콤팩트인 것이다.
\end{enumerate}
\end{proposition}

(vi)는 사상의 준콤팩트성이 바탕공간 사이의 대응하는 연속사상에만
의존하므로 명백하다. 마찬가지로 $X$의 바탕공간이 국소 뇌터라고 가정한
경우에 해당하는 (v)를 증명하자. $h=g\circ f$라 놓고, $U$를 $Y$의
준콤팩트 열린집합이라 하자. $g(U)$는 $Z$에서 준콤팩트이지만 반드시
열린집합인 것은 아니므로, 유한개의 준콤팩트 열린집합 $V_j$의 합집합에
포함된다(\hyperref[I.2.1.3-ko]{2.1.3}). 따라서 $f^{-1}(U)$는
$h^{-1}(V_j)$들의 합집합에 포함된다. 이들은 $X$의 준콤팩트
부분공간들이므로 뇌터 부분공간들이다. 그러므로
(\hyperref[0.2.2.3-ko]{0, 2.2.3})에 따라 $f^{-1}(U)$는 뇌터 공간이고,
\emph{따라서 특히} 준콤팩트이다.

나머지 주장들을 증명하기 위해서는 (i), (ii), (iii)를 증명하면 충분하다
(\hyperref[I.5.5.12-ko]{5.5.12}). 그런데 (ii)는 명백하고, (i)는 $Y$의
바탕공간이 국소 뇌터이거나 $X$의 바탕공간이 뇌터인 경우에는
(\hyperref[I.6.3.5-ko]{6.3.5})에서 따르며 닫힌 몰입 사상의 경우에는
명백하다. (iii)를 증명하기 위하여,
\oldpage[I]{154}
$S=Y$인 경우로 환원할 수 있다(\hyperref[I.3.3.11-ko]{3.3.11}).
$f'=f_{(S')}$라 놓고, $U'$를 $S'$의 준콤팩트 열린집합이라 하자. 모든
$s'\in U'$에 대하여 $T$를 $S$에서 $g(s')$의 아핀 열린 근방이라 하고,
$W$를 $U'\cap g^{-1}(T)$에 포함되는 $s'$의 아핀 열린 근방이라 하자.
$f'^{-1}(W)$가 준콤팩트임을 보이면 충분하다. 다시 말해 $S$와 $S'$가
아핀인 경우에 $X\times_S S'$의 바탕공간이 준콤팩트임을 보이는 문제로
환원된다. 그런데 이때 $X$는 가정에 따라 유한개의 아핀 열린집합 $V_j$의
합집합이고, $X\times_S S'$의 바탕공간은 아핀 스킴
$V_j\times_S S'$들의 바탕공간의 합집합이므로
(\hyperref[I.3.2.2-ko]{3.2.2}와 \hyperref[I.3.2.7-ko]{3.2.7})에 따라 명제의
증명이 끝난다.

또한 $X=X'\amalg X''$가 두 준스킴의 합이면, 사상 $f:X\to Y$가
준콤팩트이기 위한 필요충분조건은 $f$를 $X'$와 $X''$에 제한한 사상들이
준콤팩트인 것임을 주목하자.

\begin{proposition}[6.6.5]
\label{I.6.6.5-ko}
$f:X\to Y$를 준콤팩트 사상이라 하자. $f$가 우세 사상이기 위한
필요충분조건은 $Y$의 모든 기약 성분의 일반점 $y$에 대하여
$f^{-1}(y)$가 $X$의 어떤 기약 성분의 일반점을 포함하는 것이다.
\end{proposition}

이 조건이 충분함은($f$가 준콤팩트 사상이라고 가정하지 않아도) 명백하다.
필요성을 보이기 위하여 $y$의 아핀 열린 근방 $U$를 생각하자.
$f^{-1}(U)$는 준콤팩트이므로 아핀 열린집합 $V_i$들의 \emph{유한}
합집합이고, $f$가 우세 사상이라는 가정에 따라 $y$는 $U$에서 어떤
$f(V_i)$의 폐포에 속한다. $X$와 $Y$가 축소라고 해도 됨은 명백하다.
$V_i$의 기약 성분의 $X$에서의 폐포는 $X$의 기약 성분이므로
(\hyperref[0.2.1.6-ko]{0, 2.1.6}), $X$를 $V_i$로, $Y$를
$\overline{f(V_i)}\cap U$를 바탕공간으로 갖는 $U$의 축소 닫힌
부분준스킴으로 바꿀 수 있다
(\hyperref[I.5.2.1-ko]{5.2.1}). 따라서 $X=\operatorname{Spec}(A)$와
$Y=\operatorname{Spec}(B)$가 아핀이고 축소인 경우의 명제를 증명하면 된다.
$f$가 우세 사상이므로 $B$는 $A$의 부분환이고
(\hyperref[I.1.2.7-ko]{1.2.7}), 이제 명제는 $B$의 모든 극소 소 아이디얼이
$B$와 $A$의 어떤 극소 소 아이디얼의 교집합이라는 사실에서 따른다
(\hyperref[0.1.5.8-ko]{0, 1.5.8}).

\begin{proposition}[6.6.6]
\label{I.6.6.6-ko}
\begin{enumerate}
  \item[\rm(i)] 모든 국소 몰입 사상은 국소 유한형이다.
  \item[\rm(ii)] 두 사상 $f:X\to Y$, $g:Y\to Z$가 국소 유한형이면
    $g\circ f$도 국소 유한형이다.
  \item[\rm(iii)] $f:X\to Y$가 국소 유한형 $S$-사상이면, 밑준스킴의
    모든 확장 $S'\to S$에 대하여
    $f_{(S')}:X_{(S')}\to Y_{(S')}$도 국소 유한형이다.
  \item[\rm(iv)] $f:X\to X'$와 $g:Y\to Y'$가 두 국소 유한형
    $S$-사상이면 $f\times_S g$도 국소 유한형이다.
  \item[\rm(v)] 두 사상의 합성 $g\circ f$가 국소 유한형이면 $f$도
    국소 유한형이다.
  \item[\rm(vi)] 사상 $f$가 국소 유한형이면 $f_{\mathrm{red}}$도
    국소 유한형이다.
\end{enumerate}
\end{proposition}

(\hyperref[I.5.5.12-ko]{5.5.12})에 따라 (i), (ii), (iii)를 증명하면
충분하다. $j:X\to Y$가 국소 몰입 사상이면, 모든 $x\in X$에 대하여
$Y$에서 $j(x)$의 열린 근방 $V$와 $X$에서 $x$의 열린 근방 $U$가
존재하여 $j$를 $U$에 제한한 사상은 닫힌 몰입 사상 $U\to V$이다
(\hyperref[I.4.5.1-ko]{4.5.1}). 따라서 이 제한은 유한형이다. (ii)를
증명하기 위하여 점 $x\in X$를 생각하자. 가정에 따라 $g(f(x))$의 열린
근방 $W$와 $f(x)$의 열린 근방 $V$가 존재하여 $g(V)\subset W$이고
$V$는 $W$ 위의 유한형이다. 더 나아가 $f^{-1}(V)$는 $V$ 위의 국소
유한형이므로(\hyperref[I.6.6.2-ko]{6.6.2}), $x$의 열린 근방 $U$가
\oldpage[I]{155}
$f^{-1}(V)$에 포함되고 $V$ 위의 유한형이 되게 할 수 있다. 따라서
$g(f(U))\subset W$이고 $U$는 $W$ 위의 유한형이다
(\hyperref[I.6.3.4-ko]{6.3.4, (ii)}). 끝으로 (iii)를 증명하기 위해서는
$Y=S$인 경우로 환원할 수 있다(\hyperref[I.3.3.11-ko]{3.3.11}). 모든
$x'\in X'=X_{(S')}$에 대하여 $x$를 $X$에서 $x'$의 상, $s$를 $S$에서
$x$의 상, $T$를 $s$의 열린 근방, $T'$를 $S'$에서 $T$의 역상, $U$를
$x$의 열린 근방으로서 그 상이 $T$에 포함되고 $T$ 위의 유한형인 것으로
잡자. 그러면 $U\times_S T'=U\times_T T'$는 $x'$의 열린 근방이고
(\hyperref[I.3.2.7-ko]{3.2.7}), $T'$ 위의 유한형이다
(\hyperref[I.6.3.4-ko]{6.3.4, (iv)}).

\begin{corollary}[6.6.7]
\label{I.6.6.7-ko}
$X$, $Y$를 $S$ 위의 국소 유한형인 두 $S$-준스킴이라 하자. $S$가 국소
뇌터이면 $X\times_S Y$는 국소 뇌터이다.
\end{corollary}

실제로 $X$는 $S$ 위의 국소 유한형이므로 국소 뇌터이고,
$X\times_S Y$는 $X$ 위의 국소 유한형이므로 역시 국소 뇌터이다.

\begin{remark}[6.6.8]
\label{I.6.6.8-ko}
명제 (\hyperref[I.6.3.10-ko]{6.3.10})과 그 증명은 사상 $f$가 국소
유한형이라고만 가정하는 경우로 곧바로 확장된다. 마찬가지로 명제
(\hyperref[I.6.4.2-ko]{6.4.2})와 (\hyperref[I.6.4.9-ko]{6.4.9})는 그
명제들에 나타나는 준스킴 $X$, $Y$가 체 $K$ 위의 국소 유한형이라고만
가정할 때에도 성립한다.
\end{remark}
